% !TeX root = ../example.tex
\section{媒体与收尾}
\sectioncoverpage[assets/sigs_night]{本节演示视频链接、参考文献与致谢。}
\subsection{视频链接}

\begin{frame}{视频链接示例}
      \centering
      \href{run:../assets/presentation_demo.mp4}{% 原创 CC0 动画；路径相对于 build/ 中的 PDF
            \XeTeXLinkBox{\fbox{\parbox[c][.30\textheight][c]{0.65\textwidth}{\centering 点击打开示例视频}}}%
      }
      \par\medskip
      {\small 原创动画：大纲 \(\rightarrow\) 幻灯片 \(\rightarrow\) 备注（6 秒，无音轨，CC0）。}
\end{frame}

\subsection{参考文献与致谢}

\begin{frame}[fragile]{参考文献与致谢用法}
      \begin{description}[汇总单页]
            \item[汇总单页] \texttt{\string\referencespage[shrink][\string\small]} 可指定字号并缩放到一页。
            \item[汇总续页] \texttt{\string\referencespage[break][\string\footnotesize]} 内容过多时自动续页。
            \item[双栏汇总] 命令末尾加 \texttt{[twocolumn]}；省略时默认为单栏。
            \item[致谢页面] \texttt{\string\backmatter} 会同步显示封面的标题、副标题、课程/课题组、作者、指导教师和日期；加入 \texttt{[notitle]} 可隐藏主标题。
      \end{description}
\end{frame}

\subsection{收尾说明}
\begin{frame}{收尾页选项}
  \begin{itemize}
    \item 参考文献页可以选择字号、续页和栏数。
    \item 致谢页可以保留或隐藏主标题。
  \end{itemize}
\end{frame}

\subsection{参考文献汇总}
\referencespage[break][\small][onecolumn] % 第三个参数改为 [twocolumn] 可使用双栏

\subsection{致谢}
\backmatter     % 不显示主标题时改为 \backmatter[notitle]
